% TeX root = ./main.tex

% First argument to \section is the title that will go in the table of contents. Second argument is the title that will be printed on the page.

%
% Structures
%

\section[Structures algébriques]{Structures algébriques}
\subsection{Axiomes superflus ?}
\textit{Exercice donné à Louis Marshal. Remarque : un élément $x$ dans un magma $E$ est dit régulier s'il vérifie} \[\forall y,z\in E,\ (xy=xz\implique y=z\quad\text{et}\quad yx=zx\implique y=z)\]
\begin{exo}
 	Soit $E$ un magma fini. Montrer que si tout élément de $E$ est régulier alors $E$ est un groupe. Que dire si $E$ est infini ?
\end{exo}

\begin{sol}
	[À rédiger]
\end{sol}

%
% Algèbre linéaire 
%

\newpage
\section[Algèbre linéaire]{Algèbre linéaire}

\subsection{Formes linéaires $\varphi$ sur $\mathcal M_n(\R)$ vérifiant $\varphi(MM')=\varphi(M'M)$}
\begin{exo}
Soit $\varphi : \mathcal M_n(\R)\to \R$ une forme linéaire vérifiant
\[\forall M,M'\in \mathcal M_n(\R),\ \varphi(MM')=\varphi(M'M)\quad \text{ et }\quad \varphi(I_n)=n\]
On pose $A=\lbrace MM'-M'M,\ M,M'\in \mathcal M_n(\R)\rbrace$. 
\begin{enumerate}
	\item Montrer que $\text{Vect}(A)=\lbrace M\in \mathcal M_n(\R)\ |\ \text{tr}(M)=0\rbrace$.
	\item En déduire que $\varphi=\text{tr}$.
\end{enumerate}
\end{exo}
\begin{sol}
\begin{enumerate}
	\item $A\subset \text{Vect}(A)$, donc par définition, $\forall i,j,k,l\in\inl{1}{n},\ E_{ij}E_{kl}-E_{kl}E_{ij}=\delta_{jk}E_{il}-\delta_{li}E_{kj}\in A$. 

	Il s'ensuit que
	\[\forall i,l\in\inl1n,\ i\neq l\implique E_{il}\in \text{Vect}(A)\quad \text{et}\quad \forall i\in\inl{2}{n},\ E_{ii}-E_{11}\in \text{Vect}(A)\]
	On pose $\mathcal F=(E_{ij})_{i\neq j}\cup(E_{ii}-E_{11})_{i\in\inl2n}$. Cette famille est clairement libre, et contient $n^2-1$ vecteurs de $\text{Vect}(A)$, d'où $\text{dim}\text{Vect}(A)\geq n^2-1$. Or $\text{Vect}(A)\subset \text{Ker tr}$, et $\text{tr}$ étant une forme linéaire non nulle, on a $\text{dim } \text{Ker tr}=n^2-1$, donc $\text{dim} \text{Vect}(A)\leq n^2-1$ donc $\text{dim} \text{Vect}(A)=n^2-1=\text{dim}\text{ Ker tr}$ puis $\text{Vect}(A)=\text{Ker tr}$.
	\item De même que pour $\text{tr}$, on a $\text{Vect}(A)\subset \text{Ker}\varphi$ donc $\text{Ker tr}\subset \text{Ker}\varphi$, donc il existe $\lambda\in\R$ tel que $\varphi=\lambda\text{tr}$. Mais $\varphi(I_n)=n=\text{tr}(I_n)$ donc $n=\lambda n$ puis $\lambda = 1$.
\end{enumerate}
\end{sol}


\subsection{Familles libres de matrices de rang $1$ de $\mathcal M_n(\R)$}
\begin{exo}
	Soit $(X_1,\dots,X_p)\in\R^n$ une famille libre de vecteurs de $\R^n$. Montrer que la famille $(X_1{}^tX_1,\dots,X_p{}^tX_p)$ est une famille libre de vecteurs de $\mathcal M_n(\R)$. Étudier la réciproque.
\end{exo}
\begin{sol}
Pour $i\in\inl1p$, on note $X_i=(x_i^{(1)},\dots,x_i^{(n)})$. On a alors, \[\forall i\in\inl1p,\ X_i{}^tX_i=(x_i^{(1)}X_i|\dots|x_i^{(n)}X_i)\]
Soit $(\lambda_1,\dots,\lambda_p)\in\R^p$ telle que $\sum_{i=1}^p\lambda_iX_i{}^tX_i=0$. Or, $(X_1,\dots,X_p)$ formant une famille libre, aucun des vecteurs la constituant n'est nul, d'où \[\forall i\in\inl1p,\ \exists l_i\in\inl1p,\ x_i^{(l_i)}\neq 0\]
Ainsi, pour $i\in\inl1p$, en regardant que la $l_i$ème colonne dans la somme nulle écrite plus haut, on a \[\sum_{j=1}^p\lambda_jx_j^{(l_i)}X_j=0\]
$(X_1,\dots,X_p)$ étant libre, on a $\forall j\in\inl1p,\ \lambda_jx_j^{(l_i)}=0$, en particulier $\lambda_ix_i^{(l_i)}=0$ donc $\lambda_i=0$ ($x_i^{l_i}\neq 0$). Finalement, ceci étant vrai pour tout $i\in\inl1p$, on a \[\forall i\in\inl1p,\ \lambda_i=0\]
et $(X_1{}^tX_1,\dots,X_p{}^tX_p)$ est libre dans $\mathcal M_n(\R)$ (il faut bien sûr préciser que la taille des matrices $X_i{}^tX_i$ est de $n\times n$, mais cela est bien clair).
\end{sol}

\subsection{P(A) diagonalisable et P'(A) inversible $\implique$ A inversible}
\begin{exo}
Soit $A\in\mathcal M_n(\C)$ et $P\in\C[X]$ tel que $P(A)$ est diagonalisable et $P'(A)$ est inversible. Montrer que $A$ est inversible.
\end{exo}

\begin{sol}
[À rédiger]
\end{sol}

%
% EVN
%

\newpage
\section[Espaces Vectoriels Normés]{Espaces Vectoriels Normés}
\subsection{CNS pour qu'un sous-groupe de $\C^*$ soit fermé}

\textit{Exercice donné par Duhalde à Amar en semaine 3.}
\begin{exo}
Donner des conditions nécessaires et suiffisantes sur $z\in\C$ pour que $G_z=\lbrace e^{itz},\ t\in\Z\rbrace$ soit un sous-groupe fermé de $\C^*$.
\end{exo}

\begin{sol}
On procède par analyse synthèse.


\textbf{Analyse :} Soit $z=a+ib\in\C$ tel que $G_z$ est un sous-groupe fermé de $\C$. Pour $t\in\Z$, on a $e^{tz}=e^{-bt}e^{ait}$. Montrons par l'absurde que $b=0$. Supposons $b\neq 0$, traîtons les deux cas possibles : 
\begin{itemize}
	\item Si b > 0, on prend $(t_n)\in\Z^\N$ telle que $t_n\longrightarrow +\infty$. Dans ce cas, $(e^{it_nz})$ est une suite d'éléments de $G_z$. Pour $n\in\N$, on a \[|e^{it_nz}|=|e^{-bt_n}e^{iat_n}|=e^{-bt_n}\longrightarrow 0\]
	Donc $e^{it_nz}\longrightarrow 0$ et, $G_z$ étant fermé, $0\in G_z$, ce qui est exclut.
	\item Si b < 0, on refait le même raisonnement en prenant $(t_n)\in\Z^\N$ tendant vers $-\infty$, et on a $0\in G_z$, ce qui est exclut.
\end{itemize}
Donc le seul cas possible est $b=0$.


Ainsi, $z\in\R$. Pour avoir plus d'intuition sur ce que peuvent être les solutions au problème, on peut regarder le cas où $G_z$ est fini. Si $G_z$ est fini (on écarte le cas $z=0$ qui est trivial), c'est un sous-groupe fini de $\C^*$, donc, si on note $n$ son cardinal, on a $G_z=\mathbb U_n$.

Soit alors $x\in G_z\backslash$ tel qu'il existe $t\in\Z^*$ et $k\in\Z$ tels que $x=e^{itz}=e^{\frac{2ik\pi}n}$. On a alors 
\begin{align*}
	e^{i\left(tz-\frac{2k\pi}n\right)} = 0 &\implique tz - \frac{2k\pi}n\equiv 0[2\pi]\\
										   &\implique \exists l\in\Z,\ z=\frac2t\left(l+\frac{2k}n\right)\pi\\
										   &\implique z\in\pi\Q
\end{align*}
On peut aussi remarquer que $z\in\pi\Q$ suffit pour que $G_z$ soit fini. 


\textbf{Synthèse :} Soit $z\in\R$. Si $z\in\pi\Q$, $G_z$ est fini, donc c'est un compact de $\C$ comme partie finie de $\C$, donc est un fermé de $\C$. C'est aussi un sous-groupe de $\C^*$, donc $z$ convient.


Sinon, $z\notin\pi\Q$. Montrons que dans ce cas $G_z$ n'est pas un femré de $\C$. En effet, si $G_z$ est fermé, on peut montrer que $G_z=\mathbb U$. D'abord, $z\Z+2\pi\Z$ est un sous-groupe dense de $\R$, puisque sinon, on aurait, pour un $\alpha\in\R$, $z\Z+2\pi\Z=\alpha\Z$, donc $z\in z\Z+2\pi\Z=\alpha\Z$ et $2\pi\in z\Z+2\pi\Z=\alpha\Z$ donc il existe $k,l\in\Z^*$ ($z$ et $2\pi$ sont non nuls) tels que $z=k\alpha$ et $2\pi=l\alpha$, d'où $z=\frac{2k\pi}l\in\pi\Q$ ce qui est exclut. La fonction $\phi:x\in\R\mapsto e^{ix}$ étant continue, on a que $\phi(z\Z+2\pi\Z)=G_z$ est dense dans $\phi(\R)=\mathbb U$. En passant à l'adhérence, on a $G_z=\mathbb U$, ce qui est exclut, parce que, par exemple, $-1\notin G_z$.
\end{sol}

\subsection{Sous-groupe compact de $GL_n(\R)$}
\textit{Exercice donné par Duhalde à Yanis en semaine 3.}
\begin{exo}
Soit $X$ une partie de $GL_n(\R)$ non vide, compacte et stable par produit. Montrer que $X$ est un sous-groupe de $GL_n(\R)$.
\end{exo}
\begin{sol}
[À rédiger]
\end{sol}

\subsection{Une norme de $\mathcal C^2([0,1],\R)$}
\textit{Exercice donné par Duhalde à PG en semaine 3.}
\begin{exo}
On considère le $\R ev$ $E=\mathcal C^2([0,1],\R)$ munit de la norme infinie. On pose $N:f\in E\mapsto ||f+2f'+f''||_{\infty}\in\R_+$. 
\begin{enumerate}
	\item Montrer que $N$ est une norme. 
	\item On pose $g=f+2f'+f''$ pour $f\in\E$. Exprimer $f$ en fonction de $g$. 
	\item $N$ et $||\cdot||_{\infty}$ sont-elles équivalentes ?
\end{enumerate}
\end{exo}
\begin{sol}
[À rédiger]
\end{sol}

\subsection{Un oral CCINP}
\textit{Exercice donné à Shems en semaine 3.}
\begin{exo}
On munit $\R[X]$ de la norme $||\cdot||_{\infty}$ définie par $\left|\left|\sum_{k=0}^{+\infty}a_kX^k\right|\right|=\text{max}\lbrace a_k,\ k\in\N\rbrace$.
\begin{enumerate}
	\item Montrer que $\mathcal U$ l'ensemble des polynômes unitaires de $\R[X]$ est fermé.
	\item Soit $Q\in\mathcal U$ non constant, on note $p=\text{deg}Q$. Montrer que 
	\[Q\text{ est scindé sur }\R\ssi \forall z\in\C,\ |Q(z)|\geq |\text{Im}(z)|^p\]
	\item Montrer que $\mathcal S$, l'ensemble des polynômes unitaires et scindés de $\R[X]$ est un fermé.
\end{enumerate}
\end{exo}


\begin{sol}
[À rédiger]
\end{sol}